\chapter{A Concrete Naming Certificate for $\RegularCauchyTailPairingUp$}
\label{ch:concrete-instances-regular-cauchy-tail-pairing-namecert}
\origin{human}

Following Bishop--Bridges constructive real analysis, the $\RegularCauchyTailPairingUp$ packet records the finite paired-tail surface by which two regular rational Cauchy names are synchronized before a diagonal or sequential-completeness consumer reaches a displayed $\RealUp$ seal. It sits over the finite observation interface of \autoref{ch:concrete-instances-streamname-namecert}, the regular rational readback of \autoref{ch:concrete-instances-regseqrat-namecert}, the dyadic tolerance ledger of \autoref{ch:concrete-instances-dyadicratcore-namecert}, the sequential-completeness handoff of \autoref{ch:concrete-instances-regular-cauchy-sequential-completeness-namecert}, and the terminal real-name boundary of \autoref{ch:concrete-instances-real-namecert}. The packet is a finite paired-tail carrier for L10 Real and Completion routes; it is not quotient stream equality, a selected limit, countable choice, host real equality, or ambient completeness of $\RealUp$.

\begin{definition}[Regular Cauchy tail pairing carrier]
\label{def:regular-cauchy-tail-pairing-carrier}
A $\RegularCauchyTailPairingUp$ carrier is a finite $\BHist$ packet
$$
T=(S_0,S_1,R_0,R_1,D_0,D_1,W,U,E,H,C,P,N)
$$
with the following displayed rows:
\begin{enumerate}
\item $S_0$ and $S_1$ are the two $\StreamNameUp$ finite-window rows opened by the paired-tail request.
\item $R_0$ and $R_1$ are the corresponding $\RegSeqRatUp$ readback rows obtained from those windows.
\item $D_0$ and $D_1$ are the two $\DyadicRatCoreUp$ tolerance ledgers, recording the requested radius, the local regularity budgets, and the remaining comparison slack for each side of the pair.
\item $W$ is the paired-tail synchronization row, recording one displayed common tail window for the two readbacks and no quotient identification of the underlying streams.
\item $U$ is the $\RegularCauchySequentialCompletenessUp$ handoff row, read only after the two windows, two readbacks, two dyadic ledgers, and paired-tail window have been displayed.
\item $E$ is the terminal $\RealUp$ seal row reached from the same paired route.
\item $H$ records componentwise $\hsame$ transport for windows, readbacks, dyadic ledgers, paired-tail synchronization, sequential-completeness handoff, Real seal, replay, provenance, and local naming.
\item $C$ records displayed $\Cont$ replay from a diagonal or sequential-completeness request through $S_0,S_1,R_0,R_1,D_0,D_1,W,U,E,H,P,N$.
\item $P$ is the $\Pkg$ provenance over $\RegularCauchyTailPairingUp$, $\StreamNameUp$, $\RegSeqRatUp$, $\DyadicRatCoreUp$, $\RegularCauchySequentialCompletenessUp$, $\RealUp$, $\BHist$, $\hsame$, $\Cont$, and $\NameCert$.
\item $N$ is the local $\NameCert_{\RegularCauchyTailPairingUp}$ row.
\end{enumerate}
The carrier compares accepted packets only through the displayed rows above. It has no coordinate for quotient stream equality, selected host limits, countable choice, trichotomy, host real equality, or ambient $\RealUp$ completeness.
\end{definition}

\begin{theorem}[Regular Cauchy tail pairing NameCert obligations]
\label{thm:regular-cauchy-tail-pairing-namecert-obligations}
For every accepted $\RegularCauchyTailPairingUp$ carrier
$$
T=(S_0,S_1,R_0,R_1,D_0,D_1,W,U,E,H,C,P,N),
$$
the local $\NameCert_{\RegularCauchyTailPairingUp}$ surface consists exactly of carrier admission for the two StreamName windows, two RegSeqRat readbacks, two DyadicRatCore tolerance ledgers, paired-tail synchronization row, RegularCauchySequentialCompleteness handoff, Real seal, componentwise transport, displayed replay, package provenance, and local naming; synchronization exactness for the common tail window $W$; classifier stability through $H$; replay stability through $C$; and ledger non-escape through $P$ and $N$.
\end{theorem}
\begin{proof}
Carrier admission is projection from \autoref{def:regular-cauchy-tail-pairing-carrier}. The accepted tuple exposes exactly the two windows, two readbacks, two dyadic ledgers, paired-tail row, sequential-completeness handoff, Real seal, and structural rows.

Synchronization exactness is the row $W$. A diagonal or sequential-completeness consumer must first read the two StreamName windows and their RegSeqRat readbacks, then compare them under the two dyadic ledgers before the common tail window is public.

The handoff rows are downstream of that paired window. The row $U$ consumes the same displayed route as a $\RegularCauchySequentialCompletenessUp$ handoff, and the row $E$ is the terminal $\RealUp$ seal reached only after $S_0,S_1,R_0,R_1,D_0,D_1,W$ have been exposed.

Finally, $H,C,P,N$ transport, replay, source, and name the same finite packet. Since no excluded quotient, selection, choice, trichotomy, host-equality, or ambient-completeness coordinate occurs in the carrier, the listed clauses exhaust the local NameCert surface.
\end{proof}

\begin{theorem}[Regular Cauchy paired-tail synchronization]
\label{thm:regular-cauchy-tail-pairing-synchronization}
Let
$$
T=(S_0,S_1,R_0,R_1,D_0,D_1,W,U,E,H,C,P,N)
$$
be an accepted $\RegularCauchyTailPairingUp$ carrier. Any diagonal or sequential-completeness consumer of $T$ must synchronize the two tails by reading $S_0,S_1$, then $R_0,R_1$, then $D_0,D_1$, and only then the common paired-tail window $W$. The synchronization is transported componentwise by $H$ and replayed by $C$ before $U$ or $E$ is consumed.
\end{theorem}
\begin{proof}
Apply the obligation theorem \autoref{thm:regular-cauchy-tail-pairing-namecert-obligations}. It orders the public rows so that the two finite windows and readbacks are visible before either dyadic tolerance ledger is used.

The paired-tail row $W$ is therefore not a quotient or selected representative. It is a displayed common window obtained after both regular rational readbacks have been inspected under their two dyadic budgets.

Componentwise transport and replay preserve that order. The rows $H$ and $C$ can move and repeat the displayed packet, but they cannot make the sequential-completeness handoff $U$ or Real seal $E$ visible before the paired-tail synchronization row has been read.
\end{proof}

\begin{theorem}[Regular Cauchy tail pairing RealUp seal handoff]
\label{thm:regular-cauchy-tail-pairing-realup-seal-handoff}
Every $\RealUp$ or completion-facing consumer of an accepted $\RegularCauchyTailPairingUp$ carrier factors through the finite route
$$
\StreamNameUp,\ \RegSeqRatUp,\ \DyadicRatCoreUp,\ \RegularCauchySequentialCompletenessUp,\ \RealUp .
$$
More explicitly, the consumer must display the two finite windows, two regular rational readbacks, two dyadic tolerance ledgers, one paired-tail synchronization row, one sequential-completeness handoff, and one terminal Real seal before any Real-facing output is named.
\end{theorem}
\begin{proof}
Start with the paired-tail synchronization theorem \autoref{thm:regular-cauchy-tail-pairing-synchronization}. It forces every accepted consumer to read $S_0,S_1,R_0,R_1,D_0,D_1,W$ in order.

The sequential-completeness handoff is then row-local. The row $U$ cites the existing $\RegularCauchySequentialCompletenessUp$ surface of \autoref{ch:concrete-instances-regular-cauchy-sequential-completeness-namecert} and consumes the same paired route rather than selecting a host limit.

Only after that finite handoff does the Real seal row $E$ become public. The structural rows $H,C,P,N$ preserve and name the route, so the handoff exports no quotient stream equality, selected limit, countable-choice witness, host real equality, or ambient $\RealUp$ completeness.
\end{proof}

\falsifiablePrediction{Every accepted $\RegularCauchyTailPairingUp$ read displays two $\StreamNameUp$ windows, two $\RegSeqRatUp$ readbacks, two $\DyadicRatCoreUp$ tolerance ledgers, one paired-tail synchronization row, one $\RegularCauchySequentialCompletenessUp$ handoff, one $\RealUp$ seal, componentwise $\hsame$ transport, displayed $\Cont$ replay, $\Pkg$ provenance, and local $\NameCert$ evidence before a diagonal or sequential-completeness consumer reaches a Real-facing output.}

\independenceWitness{streamname, regseqrat, dyadicratcore, regular-cauchy-sequential-completeness, real: $\RegularCauchyTailPairingUp$ is not bijective to any listed sibling because it jointly carries two windows, two regular readbacks, two dyadic ledgers, one paired-tail synchronization row, one sequential-completeness handoff, one Real seal, and the structural transport, replay, provenance, and local naming rows.}

\closureat{RegularCauchyTailPairingUp}{seedStr}

\begin{closurestatus}{\RegularCauchyTailPairingUp}
  \constructivestory{A $\RegularCauchyTailPairingUp$ packet is generated from two StreamName windows, two RegSeqRat readbacks, two DyadicRatCore tolerance ledgers, one paired-tail synchronization row, one RegularCauchySequentialCompleteness handoff, one RealUp seal, componentwise hsame transport, displayed Cont replay, Pkg provenance, and a local NameCert row.}
  \theoryclosure{\seedClosure}
  \scopeclosed{\autoref{def:regular-cauchy-tail-pairing-carrier}, \autoref{thm:regular-cauchy-tail-pairing-namecert-obligations}, \autoref{thm:regular-cauchy-tail-pairing-synchronization}, and \autoref{thm:regular-cauchy-tail-pairing-realup-seal-handoff} seed the finite paired-tail carrier used by diagonal and sequential-completeness consumers before the displayed $\RealUp$ seal is read.}
  \formalstatus{\unformalizedV}
  \bridgestatus{none}
  \notclaimed{This chapter does not close quotient stream equality, selected host limits, countable choice, trichotomy, host real equality, ambient Real completeness, Lean verification, or bridge checking.}
  \upgradepath{Obligation closure requires checked carrier admission, paired-tail synchronization, classifier transport, replay stability, sequential-completeness handoff, Real seal handoff, provenance exhaustion, and local naming for BEDC.Derived.RegularCauchyTailPairingUp.}
\end{closurestatus}
